\chapter{Classification Algorithms}\doublespacing
\label{Chapter2}
\lhead{Chapter II.}

This chapter reviews the eight classifiers that make up the benchmark
of the thesis. Two of them --- Logistic Regression and $k$-Nearest
Neighbours --- are classical statistical methods. Five --- Decision
Tree, Support Vector Machine, Random Forest, Gradient Boosting and
XGBoost --- are machine-learning models. The eighth, the Stacked
Ensemble, combines the others through a meta-learner; it is
introduced formally in Section~\ref{sec:stacking}.

Throughout the chapter let
$D=\{(\mathbf{x}_i,y_i)\}_{i=1}^{N}$ denote the training set, with
covariate vector $\mathbf{x}_i\in\mathbb{R}^{p}$ and binary label
$y_i\in\{0,1\}$, where $y=1$ encodes default.

\section{Logistic Regression}
\label{sec:lr}

Logistic regression is a generalised linear model with a binomial
response and a logit link function. It models the conditional
probability of default as
\begin{equation}
P(y_i=1\mid \mathbf{x}_i)
\;=\;
\sigma\!\left(\beta_0 + \boldsymbol{\beta}^{T}\mathbf{x}_i\right),
\qquad
\sigma(t)=\frac{1}{1+e^{-t}},
\tag{2.1}
\label{eq:lr}
\end{equation}
where $\beta_0\in\mathbb{R}$ is the intercept and
$\boldsymbol{\beta}\in\mathbb{R}^{p}$ is the vector of regression
coefficients. Estimation is by maximum likelihood: the parameters
$(\beta_0,\boldsymbol{\beta})$ are chosen to maximise the
log-likelihood
\begin{equation}
\ell(\beta_0,\boldsymbol{\beta})
=
\sum_{i=1}^{N}
\bigl\{
y_i\log\sigma(\eta_i)
+
(1-y_i)\log\bigl(1-\sigma(\eta_i)\bigr)
\bigr\},
\quad
\eta_i = \beta_0 + \boldsymbol{\beta}^{T}\mathbf{x}_i.
\tag{2.2}
\label{eq:lr-ll}
\end{equation}
Logistic regression imposes a linear decision boundary in the original
feature space; its strength is interpretability (each coefficient
$\beta_j$ has the log-odds interpretation), and its weakness is its
inability to capture non-linear interactions without explicit feature
engineering. In this thesis it serves as the classical statistical
baseline.

\section{\texorpdfstring{$k$-Nearest Neighbours}{k-Nearest Neighbours}}
\label{sec:knn}

The $k$-nearest-neighbours rule is the simplest non-parametric
classifier: it stores the entire training set $D$ and, at prediction
time for a new applicant $\mathbf{x}^{*}$, retrieves the $k$ training
points closest to $\mathbf{x}^{*}$ under a chosen distance and assigns
the majority class among them.

Let $d(\mathbf{x}^{*},\mathbf{x}_i)$ denote the Euclidean distance
between $\mathbf{x}^{*}$ and $\mathbf{x}_i$, and let
$\mathcal{N}_k(\mathbf{x}^{*})$ be the set of indices of the $k$
training points with the smallest $d(\mathbf{x}^{*},\cdot)$. Then
\begin{equation}
\hat{y}(\mathbf{x}^{*})
=
\mathrm{mode}
\bigl\{
y_i \,:\, i\in\mathcal{N}_k(\mathbf{x}^{*})
\bigr\}.
\tag{2.3}
\label{eq:knn}
\end{equation}
$k$-NN is sensitive to the scale of the features --- the Euclidean
distance is dominated by attributes with large numerical range ---
so features are standardised before training. The hyperparameter
$k$ trades off variance ($k$ small) against bias ($k$ large); we use
$k=5$ throughout, which is the common default in scikit-learn.

\section{Decision Tree}
\label{sec:dt}

A decision tree partitions the feature space by a sequence of axis-
aligned binary splits, fitted greedily by maximising an impurity
reduction criterion at each node. Let
$T$ be a tree, let $\mathcal{R}_m\subset\mathbb{R}^{p}$ be the
hyperrectangle associated with the $m$-th leaf, and let $p_{m,c}$
denote the empirical class proportion in that leaf:
\begin{equation}
p_{m,c}
=
\frac{1}{N_m}
\sum_{i:\,\mathbf{x}_i\in\mathcal{R}_m}
\mathbb{I}(y_i = c),
\qquad c\in\{0,1\},
\tag{2.4}
\end{equation}
where $N_m$ is the number of training points falling into leaf $m$.
The Gini impurity of leaf $m$ is
\begin{equation}
G(m)
=
\sum_{c=0}^{1}
p_{m,c}(1-p_{m,c})
=
1 - p_{m,0}^{2} - p_{m,1}^{2}.
\tag{2.5}
\label{eq:gini}
\end{equation}
Splits are chosen to minimise the weighted sum of child impurities,
i.e.\ to maximise $G(\text{parent})-(N_L G(L)+N_R G(R))/N$ over all
possible split positions on each feature, where $L$ and $R$ are the
left and right children of the current node. The tree is grown
until a minimum-leaf or maximum-depth stopping criterion is satisfied;
in this thesis we keep the scikit-learn defaults except for setting
a fixed random seed.

\section{Support Vector Machine}
\label{sec:svm}

The Support Vector Machine finds the linear
hyperplane that maximises the margin between the two classes in a
feature space possibly transformed by a kernel. For a feature map
$\phi:\mathbb{R}^{p}\to\mathcal{H}$ inducing a kernel
$k(\mathbf{x},\mathbf{x}')=\langle\phi(\mathbf{x}),\phi(\mathbf{x}')\rangle$,
the soft-margin SVM solves
\begin{equation}
\min_{\mathbf{w},b,\boldsymbol{\xi}}
\tfrac{1}{2}\lVert\mathbf{w}\rVert^{2}
+ C\sum_{i=1}^{N}\xi_i
\quad\text{s.t.}\quad
\tilde y_i\bigl(\mathbf{w}^{T}\phi(\mathbf{x}_i)+b\bigr)\geq 1-\xi_i,
\quad \xi_i\geq 0,
\tag{2.6}
\label{eq:svm}
\end{equation}
where $\tilde y_i=2y_i-1\in\{-1,+1\}$, $C>0$ controls the trade-off
between margin width and misclassification, and the slack variables
$\xi_i$ allow individual training points to lie inside the margin.
In this thesis we use the Gaussian (radial-basis-function) kernel
\begin{equation}
k(\mathbf{x},\mathbf{x}')
=\exp\!\left(-\gamma\lVert\mathbf{x}-\mathbf{x}'\rVert^{2}\right),
\tag{2.7}
\end{equation}
which allows non-linear decision boundaries while retaining the
maximum-margin formulation.

\section{Random Forest}
\label{sec:rf}

A Random Forest \citep{breiman2001random} is an ensemble of $B$
decision trees grown on bootstrap samples of the training data, with
the additional randomisation that at each split only a random subset
of $m\leq p$ features is considered for the splitting decision. Let
$\hat T^{(b)}(\mathbf{x})$ denote the prediction of the $b$-th tree.
The forest's prediction is the majority vote
\begin{equation}
\hat y_{\mathrm{RF}}(\mathbf{x})
=
\mathrm{mode}
\bigl\{\hat T^{(1)}(\mathbf{x}),\ldots,\hat T^{(B)}(\mathbf{x})\bigr\}.
\tag{2.8}
\label{eq:rf}
\end{equation}
The bootstrap sampling and random feature subsetting decorrelate the
individual trees so that their averaged variance is smaller than the
variance of any single tree. In our experiments $B=300$ and
$m=\sqrt{p}$.

\section{Gradient Boosting}
\label{sec:gb}

Gradient Boosting also builds an ensemble
of trees, but the trees are added sequentially, each fitting the
\emph{pseudo-residuals} of the previous ensemble. Starting from a
constant predictor $F_0(\mathbf{x})$, the $m$-th boosting iteration
solves
\begin{equation}
h_m
=
\arg\min_{h}
\sum_{i=1}^{N}
\bigl(
r_{i,m} - h(\mathbf{x}_i)
\bigr)^{2},
\quad
r_{i,m}
=
-\!\left.
\frac{\partial L(y_i,F(\mathbf{x}_i))}{\partial F(\mathbf{x}_i)}
\right|_{F=F_{m-1}},
\tag{2.9}
\label{eq:gb}
\end{equation}
where $L$ is the log-loss for binary classification and $h$ ranges
over depth-limited regression trees. The new predictor is
$F_m = F_{m-1} + \nu h_m$, with shrinkage factor $\nu\in(0,1]$.
After $M$ boosting iterations the final classifier converts
$F_M(\mathbf{x})$ to a probability via the logistic transformation
$\sigma(F_M(\mathbf{x}))$.

\section{XGBoost}
\label{sec:xgb}

XGBoost \citep{chen2016xgboost} is a regularised, second-order
implementation of gradient boosting that has become the most popular
single-classifier choice for tabular benchmarks. Its objective at
iteration $m$ includes both the loss and an explicit complexity
penalty:
\begin{equation}
\mathcal{L}^{(m)}
=
\sum_{i=1}^{N}
L\!\left(y_i,F_{m-1}(\mathbf{x}_i)+h_m(\mathbf{x}_i)\right)
+
\Omega(h_m),
\quad
\Omega(h)=\gamma T + \tfrac{1}{2}\lambda\lVert\mathbf{w}\rVert^{2},
\tag{2.10}
\label{eq:xgb}
\end{equation}
where $T$ is the number of leaves of $h_m$, $\mathbf{w}$ is the vector
of leaf scores, and $\gamma,\lambda\geq 0$ are regularisation
hyperparameters. XGBoost uses a second-order Taylor expansion of $L$
to compute the leaf scores and split gains in closed form, which
yields both faster training and tighter regularisation than the
first-order gradient-boosting recursion.

\section{Stacked Ensemble}
\label{sec:stacking}

Stacked Generalisation, introduced by \citet{wolpert1992stacked}, is
a meta-learning method in which an additional model --- the
\emph{meta-learner} $g$ --- is trained to combine the predictions of
several \emph{base learners} $h_1,\ldots,h_M$. The principal
challenge is to avoid \emph{data leakage}: if the meta-features were
the in-sample predictions of the base learners, the meta would be
trained on values that are already too optimistic.

\subsection{Out-of-fold Meta-features}

Stacking solves the leakage problem by $k$-fold out-of-fold (OOF)
prediction. Partition the training set into $k$ folds. For each
fold $j$ and each base learner $h_m$:
\begin{enumerate}
 \item Train $h_m$ on the data of all folds except $j$.
 \item Predict on the held-out fold $j$.
\end{enumerate}
Every training point thus receives one prediction from a base learner
that did not see it during training. Stacking these predictions
across the $M$ base learners gives the \emph{meta-feature matrix}
\begin{equation}
\mathbf{z}_i
=
\bigl(
h_1^{(-)}(\mathbf{x}_i),\,
h_2^{(-)}(\mathbf{x}_i),\,
\ldots,\,
h_M^{(-)}(\mathbf{x}_i)
\bigr)^{T}
\in\mathbb{R}^{M},
\tag{2.11}
\label{eq:meta-features}
\end{equation}
where the superscript $(-)$ indicates the OOF prediction.

\subsection{Meta-learner Training}

The meta-learner $g$ is then estimated on the pairs
$\{(\mathbf{z}_i,y_i)\}_{i=1}^{N}$ by minimising a loss
$L_{\text{meta}}$:
\begin{equation}
g^{*}
=
\arg\min_{g}
\sum_{i=1}^{N}
L_{\text{meta}}\!\bigl(y_i,\,g(\mathbf{z}_i)\bigr).
\tag{2.12}
\label{eq:meta-train}
\end{equation}
At prediction time, each base learner is refit on the \emph{whole}
training set, the new applicant $\mathbf{x}^{*}$ is pushed through
all the base learners to obtain a fresh meta-feature vector
$\mathbf{z}^{*}$, and the stacked prediction is
$\hat y_{\mathrm{stack}}(\mathbf{x}^{*})
= g^{*}(\mathbf{z}^{*})$.

The full pipeline is summarised in
Figure~\ref{fig:stacker} (Chapter~\ref{Chapter3}), where the
specific choice of base and meta learners used in this thesis is
also reported. We defer that material to Chapter~\ref{Chapter3}
because the meta-selection rule and the GA-MaxAccLogit
candidate are themselves methodological contributions of the work.

\subsection{Why Stacking Works}

The expected error of a classifier can be decomposed into the
sum of its squared bias and its variance. Averaging the
predictions of diverse classifiers --- as bagging does --- reduces
the variance term but leaves the bias unchanged. Stacking, in
contrast, lets the meta-learner \emph{re-weight} the base learners
in a region-dependent way: when the OOF predictions reveal that
one base learner is more accurate in a certain area of feature
space, the meta-learner can lean on that base learner there and
on a different base learner elsewhere. In effect the meta-learner
can correct systematic bias of any individual base learner so
long as the bias is detectable from the OOF predictions of the
other base learners. This is the
intuition behind the empirical observation, confirmed in
Chapter~\ref{Chapter5}, that the stacked ensemble is the only
classifier in the benchmark that is never the weak link on any
of the four datasets.
